\subsection{Проектирование функциональных требования}

Функционал программного средства – ключевая точка разработки
программного средства, составляющей которого являются
алгоритмы.
В связи с этим, программное средство {\taskNameFull} должно
обладать некоторыми алгоритмами:
\begin{itemize}
    \item алгоритм загрузки настроек приложения из реестра;
    \item алгоритм расширения узла дерева каталога;
    \item алгоритм расширения сетевого узла;
    \item алгоритм обновления строки состояния директории;
    \item алгоритм обработки файлов директории;
\end{itemize}

\subsubsection{Алгоритм загрузки настроек приложения}

Система управления настройками позволяет сохранять информацию
о пользовательских предпочтениях (тема оформления, размер
и тип шрифта) в реестре Windows и предоставляет возможность
восстановления этих параметров при следующем запуске
приложения.

Алгоритм загрузки настроек приложения из реестра представлен
на рисунке \ref{fig:settings}.

\begin{figure}[ht]
    \centering
    \includegraphics[width=0.60\linewidth]{\commonSecPathPrefix/sec_2/content/settings.jpg}
    \caption{Алгоритм загрузки настроек приложения из реестра}
    \label{fig:settings}
\end{figure}

Перед началом работы необходимо открыть ключ реестра, в
котором хранятся настройки приложения. Далее, при успешном
открытии ключа, последовательно считываются три параметра:
тема оформления, размер шрифта и имя шрифта. Для каждого
параметра выполняется проверка успешности чтения и валидация
значения. Если параметр успешно
прочитан и валиден, он применяется к текущей конфигурации.
После обработки всех параметров ключ реестра закрывается,
и вызывается функция обновления цветов темы для применения
загруженных настроек к интерфейсу приложения.

\subsubsection{Алгоритм расширения узла дерева каталога}

Система навигации по файловой системе -- это система, которая
позволяет отображать структуру каталогов в виде дерева и
обеспечивает динамическую загрузку подкаталогов.

Система навигации по файловой системе позволяет отображать
иерархическую структуру директорий, обеспечивает ленивую
загрузку содержимого и предоставляет возможность эффективной
навигации по большим файловым системам.

Алгоритм расширения узла дерева каталога представлен на рисунке \ref{fig:expand}.

\begin{figure}[ht]
    \centering
    \includegraphics[width=0.60\linewidth]{\commonSecPathPrefix/sec_2/content/expand.jpg}
    \caption{Алгоритм расширения узла дерева каталога}
    \label{fig:expand}
\end{figure}

Перед началом работы необходимо удалить все
placeholder-элементы,
которые были добавлены для обозначения возможности расширения
узла. Далее проверяется, был ли узел уже расширен ранее,
чтобы избежать повторной загрузки. Если узел ещё не расширен,
открывается поиск файлов в соответствующей директории. При
успешном открытии поиска выполняется перебор всех найденных
элементов: пропускаются системные директории,
для каждой найденной директории проверяется наличие
подкаталогов, и элемент добавляется в дерево с
соответствующей пометкой о наличии дочерних элементов.
После завершения перебора поиск закрывается, и узел
помечается как расширенный.

\subsubsection{Алгоритм расширения сетевого узла}

Система работы с сетевыми ресурсами -- это система,
которая позволяет обнаруживать и отображать сетевые
устройства и общие ресурсы в локальной сети.

Система работы с сетевыми ресурсами позволяет находить
сетевые серверы и шары, обеспечивает динамическую загрузку
списка доступных ресурсов и предоставляет возможность
навигации по сетевым директориям.

Алгоритм расширения сетевого узла представлен на рисунке \ref{fig:network}.

\begin{figure}[ht]
    \centering
    \includegraphics[width=0.60\linewidth]{\commonSecPathPrefix/sec_2/content/network.jpg}
    \caption{Алгоритм расширения сетевого узла}
    \label{fig:network}
\end{figure}

Перед началом работы необходимо удалить все
placeholder-элементы из узла сетевого устройства.
Далее проверяется, был ли узел уже расширен, чтобы избежать 
повторной загрузки сетевых ресурсов. Если узел ещё не
расширен, определяется тип сетевого пути: проверяется
наличие третьей обратной косой черты в пути. Если третья
черта не найдена, это означает, что путь указывает на
сетевой сервер, и выполняется
перебор всех доступных сетевых ресурсов на этом
сервере. Для каждого найденного ресурса извлекается его
имя, проверяется наличие подкаталогов, и ресурс добавляется
в дерево. Если третья черта найдена, это означает, что путь
указывает на конкретную шару или подкаталог, и для его
обработки вызывается стандартный алгоритм расширения
локального узла дерева.

\subsubsection{Алгоритм обновления строки состояния директории}

Система отображения статистики директории -- это система,
которая позволяет отображать информацию о количестве файлов
и папок в текущей директории.

Система отображения статистики директории позволяет
подсчитывать количество элементов различных типов,
суммировать размеры файлов и предоставляет возможность
отображения этой информации пользователю в строке состояния
интерфейса.

Алгоритм обновления строки состояния директории представлен на рисунке \ref{fig:status}.

\begin{figure}[ht]
    \centering
    \includegraphics[width=0.60\linewidth]{\commonSecPathPrefix/sec_2/content/status.jpg}
    \caption{Алгоритм обновления строки состояния директории}
    \label{fig:status}
\end{figure}

Перед началом работы необходимо проверить существование окна
строки состояния. Далее инициализируются счётчики файлов
и папок. Выполняется перебор всех элементов в текущей
директории: для каждого элемента определяется его тип. Если
элемент является директорией,
увеличивается счётчик папок. Если элемент является файлом,
увеличивается счётчик файлов, и выполняется попытка открыть
файл для получения его размера. При успешном открытии файла
его размер добавляется к общему размеру, после чего файл
закрывается. После обработки всех элементов проверяется,
есть ли в директории файлы и папки: если директория пуста,
в строку состояния устанавливается текст "Готово", иначе
формируется текст с количеством файлов и папок. Сформированный текст отправляется в
окно строки состояния для отображения пользователю.

\subsubsection{Алгоритм обработки файлов директории}

Система отображения содержимого директории -- это система,
которая позволяет получать список файлов и папок в
директории и отображать их в интерфейсе приложения.

Система отображения содержимого директории позволяет
перебирать элементы файловой системы, определять их типы,
получать системные иконки и предоставляет возможность
отображения элементов в виде списка с информацией о типе
и времени последнего изменения.

Алгоритм обработки файлов директории представлен на рисунке \ref{fig:files}.

\begin{figure}[ht]
    \centering
    \includegraphics[width=0.60\linewidth]{\commonSecPathPrefix/sec_2/content/files.jpg}
    \caption{Алгоритм обработки файлов директории}
    \label{fig:files}
\end{figure}

Перед началом работы инициализируется счётчик элементов.
Выполняется перебор всех найденных файлов и директорий через
функцию поиска файловой системы: для каждого найденного
элемента проверяется его имя. Если имя элемента равно "."
или "..", элемент пропускается.
Для остальных элементов определяется тип, создаётся объект
элемента списка с именем
и полным путём, и элемент добавляется в вектор элементов.
Далее для каждого элемента получается системная иконка,
соответствующая его типу, создаётся элемент списка интерфейса
с установкой текста, параметра и иконки. Элемент вставляется
в список интерфейса, после чего устанавливается текст
колонки "Type" и
текст колонки "Modified" с информацией о времени последнего
изменения. Счётчик элементов увеличивается, и обработка
продолжается до тех пор, пока есть следующие файлы. После
завершения перебора поиск файлов закрывается.
